% !TEX root = ../main.tex
% =====================================================================
% 3. THE DESIGN SPACE -- page budget 1.25
%
% This is the "method" section, but it is written as the axes of the
% study, not as a product pitch. Full equations for every head and every
% encoder variant are in Appendix B; here we give the two slots, the
% reduction properties (which is what makes the sweep controlled), and
% only the one head that Section 5 needs in closed form.
% =====================================================================

\section{The Design Space We Sweep}
\label{sec:design}

A model in this family is a pair \emph{encoder $+$ pooling head}. The encoder consumes
$(B,1,\nsteps)$ and returns $(B,\dmodel,\nsteps)$; the head returns the prediction $(B,\nclz)$ and
the interpretation $(B,\nclz,\nsteps)$. Every component below respects that contract, so any encoder
can be paired with any head from a configuration file alone. That is what makes the sweep of
\Cref{sec:protocol} controlled rather than a collection of separate models.

\subsection{The encoder axis: capacity at fixed structure}
\label{sec:design:enc}

We use \seanet{}, a length-preserving convolutional encoder: a stem convolution followed by
\emph{multi-scale separable residual blocks}. One unit sums three depthwise convolutions of kernel
size $\{5,11,23\}$ at a shared dilation, mixes channels with a pointwise $1\times1$ convolution, and
is wrapped in BatchNorm, ReLU, dropout and a residual connection. Dilation grows geometrically and
is capped at 16. Three properties matter for this study, and all three are constraints rather than
features. Padding is ``same'' with no stride and no pooling, so $\nsteps$ is preserved at every
stage -- required, because the interpretation has one value per step and AOPCR masks individual
steps. The dilation cap keeps each embedding tied to a bounded window, so importance maps stay
local. And splitting each convolution into a depthwise and a pointwise step
\citep{chollet2017xception, howard2017mobilenets} costs $\dmodel k + \dmodel^2$ weights instead of
$\dmodel^2 k$, which is what makes a small setting reachable at all.

The encoder is instantiated at two capacities, and \Cref{sec:cheap} turns that into the paper's
capacity control: \textbf{narrow} ($\dmodel = 64$, 4 blocks, $41$--$70$\,K parameters) and
\textbf{wide} ($\dmodel = 128$, 6 blocks, $\approx 269$\,K parameters). Seven variants sit on top of
the same block -- bottleneck, gated, input-gated, spike/trend, reconstruction-residual, and two
wrappers (multi-scale input channels, multi-scale pyramid). They are described in
\Cref{app:encoders}; for this paper they matter only as the spread of the encoder axis, which
\Cref{sec:cheap} compares against the spread of the head axis.

\subsection{The aggregation axis: seven heads in two slots}
\label{sec:design:heads}

Every head keeps the two-branch template of \Cref{eq:conj_gate} -- an attention branch deciding
\emph{where} to look and a per-step classifier deciding \emph{what} each position argues for -- so a
head is fully specified by two choices: the shape of the attention branch (slot~1) and the operator
that reduces the $\nsteps$ evidence values $g_j^k$ to one number (slot~2). \Cref{tab:heads} lays the
seven out on those two axes.

\begin{table}[t]
\caption{The seven pooling heads as two independent slots. $g_j^k$ is the evidence of step $j$ for
class $k$. ``Reduces to baseline'' names the parameter setting at which the head becomes conjunctive
pooling \emph{exactly}; this is what makes the $\kappa$ sweep of \Cref{sec:selection} a controlled
experiment, and what \Cref{sec:safety} shows is not sufficient in practice. L1--L3 are the
limitations of \Cref{sec:bg:conj}.}
\label{tab:heads}
\begin{center}
\scriptsize
\setlength{\tabcolsep}{3.5pt}
\begin{tabular}{lllcc}
\toprule
\textbf{Head} & \textbf{Slot 1: attention} & \textbf{Slot 2: reduction over time} &
\textbf{Reduces at} & \textbf{Targets} \\
\midrule
Conjunctive (baseline)    & shared gate    & mean                            & --- & --- \\
Class-wise conjunctive    & per-class gate & mean                            & $a_j^k = a_j$ & L1 \\
Softmax conjunctive       & per-class score& softmax, learnable $\tau$       & $\tau \to \infty$ & L2 \\
Adaptive class-wise       & per-class gate & soft-argmax, learnable $\beta_k$& $\beta_k \to 0$ & L3 \\
Top-$k$ conjunctive       & per-class gate & mean of the $k$ largest         & $\kappa = 1$ & L3 \\
Attention-max             & per-class gate & blend of mean and max           & $\lambda_k \to 0$ & L3 \\
Per-class gated attention & gated $\tanh\!\odot\!\sig$ & softmax             & \textbf{none} & L1, L3 \\
Dual-stream conjunctive   & gate $+$ critical query & blend of two streams   & $\lambda \to 0$ & L3 \\
\bottomrule
\end{tabular}
\end{center}
\end{table}

Two of the eight rows are adapted from histopathology MIL rather than invented here: gated attention
is due to \citet{ilse2018attention}, and the dual-stream mechanism to \citet{li2021dsmil}. Our
contribution to those two is the transfer to time series, the per-class formulation, and -- for the
dual stream -- removing its contrastive pre-training and adding the learnable blend that gives it
the reduction property in the fourth column.

\paragraph{Top-$k$ conjunctive pooling, in full.}
\Cref{sec:selection} rests on this head, so we give it here. Let $g_j^k = a_j^k \instpred_j^{\,k}$
with $a_j^k$ the per-class gate. Instead of averaging all $\nsteps$ evidence values, average only the
$k$ largest, with $k$ a fixed fraction $\kappa$ of the series length:
\begin{equation}
  k = \max\big(1, \lceil \kappa \nsteps \rceil\big),
  \qquad
  \bagpred^{\,k} = \frac{1}{k}\!\!\sum_{j \in \mathcal{S}^k_{\mathrm{top}}}\!\! g_j^k,
  \qquad
  \interp_{k,j} = g_j^k ,
  \label{eq:topk}
\end{equation}
where $\mathcal{S}^k_{\mathrm{top}}$ holds the indices of the $k$ largest $g_j^k$. Note that
$\kappa$ is a hyperparameter, not a weight: this head has \emph{exactly the same parameter count} as
class-wise conjunctive pooling.

\begin{property}[Exact reduction]
\label{prop:topk}
At $\kappa = 1$, \Cref{eq:topk} is class-wise conjunctive pooling, term for term. For $\kappa < 1$
the gradient reaches only the selected steps.
\end{property}

Property~\ref{prop:topk} is the reason \Cref{sec:selection} can attribute its result to selection
and to nothing else: the $\kappa = 1$ row of that experiment is not a different model, it is the
same model with the selection switched off.

\subsection{Training objective}
\label{sec:design:loss}

Every model is trained with cross-entropy on the bag logits plus one term that exists for the
explanation rather than for the classification:
\begin{equation}
  \mathcal{L} = \mathrm{CE}(\bagpred, y)
  \;+\; \lambda_{\mathrm{focus}} \Big({-}\textstyle\sum_{j} \bar a_j \log(\bar a_j + \epsilon)\Big),
  \label{eq:loss}
\end{equation}
with $\bar a_j$ the class-averaged attention. Attention spread evenly over the series has high
entropy, so adding entropy to a minimised loss pushes the model towards concentrated attention.
Without it, a model can be right while its importance map is flat, and a reader cannot tell where
the evidence stops. This is the second, independent selection knob used in \Cref{sec:selection}. We
set $\lambda_{\mathrm{focus}} = 0.01$ for every \seanet{} model and $0$ for the \millet{} baseline,
since the term is our addition; label smoothing is $0.13$ throughout.
